\section{The explanation and the checks}
\label{sec:checks}

\ifjof
\paragraph{Use of AI tools} Karthik Iyer designed the audit. He chose the retail question, the explanation checks and
their pass rules, the generated stores with known effects, and the measures of what a wrong lift costs in stock and
promotion decisions. Claude (Anthropic) assisted with writing and testing code for data preparation, modelling,
experiments and analysis. Iyer reviewed the results, identified implementation errors and directed corrections. He
specified the figures and how the results should be presented, and reviewed successive versions. The figures and
tables were produced from the experiment outputs using scripts developed with Claude's assistance. The author takes
responsibility for the final work.
\fi

\subsection{Notation}

A data set has series $i$ (an item in a store) observed in periods $t$. The
plan for series $i$ in period $t$ is $x_{it} = (p_{it}, \ell^{1}_{it}, \dots,
\ell^{L}_{it})$: the shelf price and $L$ lever indicators (a deal flag on
Dominick's; promotion, display, feature and a tag-only price reduction
flag on Frat; event and SNAP days on M5; promotion, event and payday days on
Favorita, which has no prices).
A model fitted on the periods before an origin forecasts the next $H$ periods
from the history and the plan for those periods, $\hat y_{it} = f_{it}(x)$.

The audit's \emph{regular price} $r_i$ is the series' highest price over the
last 12 training periods, or the planned price where the series has no price in
those periods. The plan with \emph{nothing planned}, $x^{0}$, sets every price
to $r_i$ and every lever to zero for all $H$ periods. Each model also has its
own regular price and its own plan for nothing planned, which can differ from
the audit's (\S\ref{sec:results}); that difference is what makes C1 fail under
the audit's plan. The explanation of a forecast has four parts:
\begin{align}
  b_{it} &= f_{it}(x^{0}) \qquad \text{baseline}\\
  \Lambda_{it} &= f_{it}(x) - b_{it} \qquad \text{total lift}\\
  c^{k}_{it} &= f_{it}(x) - f_{it}(x \mid \ell^{k} = 0) \qquad \text{lever } k\\
  c^{p}_{it} &= f_{it}(x) - f_{it}(x \mid p = r) \qquad \text{price}
\end{align}
and the interaction $\iota_{it} = \Lambda_{it} - c^{p}_{it} - \sum_k c^{k}_{it}$,
which is zero only when the levers act separately. All four are in units.
Switching a lever off switches it off for all $H$ periods, so the dip after a
promotion is part of that promotion's contribution.

The world model and the lift model of the \companion{} ~\cite{rwm} compute
the split themselves. Their network keeps a \emph{state} for each series: an
\emph{encoder} turns the history into the starting state, and a
\emph{transition} steps the state forward one period at a time given that
period's plan (the state is \emph{rolled} through the plan). Two output layers
read each period's state: the \emph{baseline readout} gives the forecast under
the model's own plan for nothing planned, and the \emph{lift readout} gives a
lift factor, the readout on the plan as made divided by the same readout on the
nothing-planned plan, which is exactly one when nothing is planned. The forecast
is the baseline times the lift factor; the explanation reports it in units, as
the baseline $b$ plus the lift $\Lambda = \hat y - b$. For
LightGBM~\cite{lightgbm} the same on/off quantities are computed from its
forecasts, together with TreeSHAP values $\phi^{k}_{it}$~\cite{lundberg2020}
on its lever and price inputs. Unlike Shapley values, on/off contributions
answer the merchant's question (what would sell without this lever) directly,
and are defined for any model.

Rank agreement $\rho(a, b)$ below is the Spearman correlation over item-weeks
where either $a$ or $b$ is non-zero, so that the many weeks with nothing planned
do not inflate it.

\subsection{The eight checks}

\paragraph{C1 By construction} Three violations, each relative to the largest
forecast: the lift with nothing planned, $\max|\Lambda(x^{0})|$; the movement
of the baseline with the plan, $\max|b(x) - b(x^{0})|$; and
$\max|b + \Lambda - \hat y|$, with $b$ and $\Lambda$ as the model's readouts
give them. These hold by design and are checked because a design can be broken
by an implementation.

\paragraph{C2 Randomisation} Following Adebayo et al.~\cite{adebayo2018}, the
network's parts are re-initialised in turn from the output back to the input:
lift readout, baseline readout, transition, encoder. After each step, the rank
agreement of the total lift with the original, $\rho(\Lambda, \Lambda^{(k)})$,
should fall. An explanation that survives randomisation does not depend on what
the model learned. The known-answer study shows that C2 does not test a lift
built into the model (\S\ref{sec:known}); from there on it is reported, and
counted in the main test as fixed in advance, but not used as evidence.

\paragraph{C3 Seed stability} Fits of the same configuration that differ only
in their seed should explain a forecast alike. The score is the mean over seed
pairs of $\rho(\Lambda^{s}, \Lambda^{s'})$, shown against the same agreement for
the forecasts themselves. For a final tuned model, which averages copies fitted
with different seeds, the pairs are its copies.

\paragraph{C4 Price direction} The planned price is cut by 5\%, 10\%, 20\% and
30\% in turn. For each item-week and each step to a deeper cut (0 to 5\%, 5\%
to 10\%, 10\% to 20\%, 20\% to 30\%), the step fails when the deeper cut gives
a forecast lower than the shallower one by more than 0.1\% of the forecast. The
score is the share of item-week steps that fail.

\paragraph{C5 Against what sold} On item-weeks with a lever on or a price at
least 5\% below regular, the bias $\sum \hat y / \sum y$ and the rank agreement
of the predicted lift factor $\hat y/b$ with the observed one $y/b$. Merchants
plan promotions for weeks they expect to sell well, so $y/b$ overstates what
the plan caused; this check measures agreement with the data, not with a cause.

\paragraph{C6 Against LightGBM} The rank agreement of the total lift with
LightGBM's on the same window, $\rho(\Lambda, \Lambda_{\mathrm{LGB}})$, and for
each lever $\rho(c^{k}, c^{k}_{\mathrm{LGB}})$ and $\rho(c^{k}, \phi^{k})$. Two
models fitted to the same data can agree and both be wrong; agreement is
evidence of a shared signal, not of correctness.

\paragraph{C7 Against a measured cause} For an episode $e$ in which store $s$'s
regular prices sat at a log ratio $\delta_e$ from its zone peers, the model's
stated response is
\begin{equation}
  \hat R_e = \log \frac{\sum_{t \in E} \sum_{i \in s} f_{it}(x)}
                      {\sum_{t \in E} \sum_{i \in s}
                       f_{it}(x \mid p_{s} = p_{s} e^{-\delta_e})},
\end{equation}
the change in the store's forecast category units if its prices had stayed with
its zone, over the episode's first eight weeks $E$. It is compared with the
measured response $\tau_e$ from difference in differences
(\S\ref{sec:zones}). The score fixed in advance is per episode: whether
$\hat R_e$ falls inside $\tau_e$'s 95\% interval and whether its sign agrees.
The comparison of the pooled slopes of $\hat R_e$ and $\tau_e$ on $\delta_e$
was added after the robustness results showed the per-episode intervals to be
too narrow, and is exploratory (Appendix~\ref{app:plan}).

\paragraph{C8 Plan read-back} For the world model only: a ridge probe reads
each plan input of a period (the price as a share of the regular and of the
usual price, each lever, and whether each is recorded) back from the model's
state after it has been rolled through the plan. The score is the $R^2$,
averaged over the inputs, against the same probe on the state before the
window, which has not seen the plan. Because the rolled state has received the
plan as input, a high $R^2$ shows only that the state keeps it; the gain over
the state before the window is the informative part.

\subsection{Pass thresholds and the main test}
\label{sec:test}

Pass rules for C1 to C6 were fixed in the plan before any audit result was read
(Table~\ref{tab:thresholds}); C7 and C8 have none and are reported per model.
A \emph{model version} is a configuration and a seed; a version-window is a
version scored on one rolling window. For the main question, each
version-window gets the number of checks it passes among C1 to C6, counting
only those that apply to it. Within each data set the Spearman correlation
across version-windows between that count and item-by-store WRMSSE (lower is
better) is computed, and the same for each check's continuous score. A negative
correlation means versions that pass more checks forecast better. The data sets
are pooled by a permutation test that shuffles the scores within each data set,
10{,}000 times.

\begin{table}[t]
\centering
\caption{\textbf{C1 to C6 have a pass rule fixed before any result; C7 and C8
do not.} Scores as defined in \S\ref{sec:checks}; a check that does not apply
to a model (C1 and C2 for LightGBM) is not counted for it.}
\label{tab:thresholds}
\footnotesize
\begin{tabular}{@{}lp{3.3cm}p{2.3cm}@{}}
\toprule
Check & Passes when & Applies to\\
\midrule
C1 & every violation $< 0.001$ & lift readout models\\
C2 & $\rho$ after the lift readout $< 0.5$ & neural models\\
C3 & mean seed $\rho \ge 0.7$ & configurations with $\ge 2$ seeds\\
C4 & share $< 1\%$ & all\\
C5 & bias in $[0.8, 1.25]$ and rank $\ge 0.3$ & all\\
C6 & $\rho(\Lambda, \Lambda_{\mathrm{LGB}}) \ge 0.5$ & non-LightGBM\\
C7 & no pass rule; per-episode score fixed in advance, pooled slope added after
the robustness results & Dominick's models\\
C8 & no pass rule & world model\\
\bottomrule
\end{tabular}
\end{table}
